% %%%%%%%%%%%%%%%%%%%%%%%%%%%%%%%%%%%%%%%%%%%%%%%%%%%%%%%%%%%%%%%%%%%%%%%%%%%%%%%%%
%
%	cap#.tex 
%  -----------
%
%	Neste arquivo você escreve o capítulo. Deve sempre iniciar com o ambiente
%	\chapter{Nome do Capitulo}
%
% %%%%%%%%%%%%%%%%%%%%%%%%%%%%%%%%%%%%%%%%%%%%%%%%%%%%%%%%%%%%%%%%%%%%%%%%%%%%%%%%%
%
\chapter{Cronograma}
% ----------------------------------------------------------
\begin{table}[H]
%	\caption{Parâmetros relacionados à geometria do domínio, à escavação e à instalação do revestimento do modelo axissimétrico}
	\label{cronograma}
	\centering
	\small
	\renewcommand{\arraystretch}{1.25}
	\begin{tabular}{|l|c|c|c|c|c|c|}
		\hline
		\textbf{Atividades} & \textbf{Out} & \textbf{Nov} & \textbf{Dez} & \textbf{Jan} & \textbf{Fev} & \textbf{Mar} \\ \hline	
		Revisão dos resultados preliminares                         &  &  &  &  &  &  \\ \hline
		Aprimorar os resultados do capítulo 3                       &  &  &  &  &  &  \\ \hline
		Incluir apêndices fundamentais                              &  &  &  &  &  &  \\ \hline
		Ampliar o caso de duas famílias ao caso genérico            &  &  &  &  &  &  \\ \hline
		Formular a propagação de dano contínuo                      &  &  &  &  &  &  \\ \hline
		Avaliar a propagação de dano contínuo                       &  &  &  &  &  &  \\ \hline
		Aplicações numéricas do trabalho                            &  &  &  &  &  &  \\ \hline
		Finalização do texto da dissertação                         &  &  &  \cellcolor{blue!30} & \\ \hline
		Desenvolvimento do artigo                                   &  &  &  &  \cellcolor{blue!30} & \cellcolor{blue!30} \\ \hline
		Defesa da dissertação                                       &  &  &  &  &  & \cellcolor{blue!30} \\ \hline
	\end{tabular}
\end{table}

\end{document}